% Chapter 3: Context Management and Memory
\chapter{Context Management and Memory}
\label{ch:context}

The context window is the fundamental constraint of AI-assisted development.
Every token of context spent on irrelevant information is a token unavailable
for the task at hand. Managing context is not housekeeping---it is the
highest-leverage skill for productive Claude usage.

\section{The Context Window as Your Scarcest Resource}
\label{sec:context-scarcity}

\marginnote[Official]{Use \code{/clear} between unrelated tasks. Use \code{/compact} for selective summarization. Use subagents for context-heavy research.\sourceurl{https://code.claude.com/docs/en/best-practices}}

\official{Performance degrades as context fills.}

\practitioner{In practice, this is not gradual---there is
a noticeable threshold where Claude begins to lose track of earlier instructions,
repeat itself, or ignore CLAUDE.md rules it followed perfectly 20 minutes ago.}

\subsection{Context Hygiene Practices}

\begin{description}
  \item[\code{/clear}] Wipe context and start fresh. Use between unrelated tasks.
  \item[\code{/compact}] Summarize conversation while preserving key decisions.
    Less aggressive than \code{/clear}; retains important context.
  \item[\code{/rewind}] Checkpoint and roll back to a previous state.
  \item[Subagents] Delegate context-heavy research to isolated agents.
    Results return to your session; the research context does not.
\end{description}

\subsection{Monitoring Context Usage}

\margintip{Watch the token counter. When it approaches 60--70\% of the model's context window, proactively \code{/compact} or \code{/clear}.}

\official{Use \code{/status} to see current session state: model, context usage,
and available tools.} Claude Code displays token consumption in the interface.
The custom status line (\code{statusLine} setting) can be configured to show
context usage prominently.

\practitioner{If you preach context awareness, you must observe context.
Check \code{/status} before starting any complex task.
If context is above 50\%, consider whether to compact before proceeding.}

\subsection{The Two-Failure Rule}

\practitioner{After two failed corrections on the same issue, \code{/clear}
and write a better initial prompt.}

The cost of accumulated bad context exceeds the cost of starting fresh.
Each correction attempt adds noise: the original error, your correction,
Claude's apology, and its retry. After two rounds, the context is
dominated by failure patterns rather than the actual task.

\begin{warnbox}[Context Pollution]
Three rounds of ``no, that's not what I meant'' produces $\sim$2000 tokens
of noise that actively degrades subsequent responses. \code{/clear} costs
30 seconds. Persisting costs quality for the rest of the session.
\end{warnbox}

% -------------------------------------------------------------------
\section{Session Management}
\label{sec:sessions}

\marginnote[Official]{Named sessions: \code{/rename}, \code{--continue}, \code{--resume}, \code{--from-pr}.\sourceurl{https://code.claude.com/docs/en/common-workflows}}

\official{Sessions persist across terminal restarts.} Use them for:

\begin{description}
  \item[\code{claude --continue}] Resume the most recent session.
  \item[\code{claude --resume}] Pick from a list of previous sessions.
  \item[\code{claude --from-pr 123}] Start a session with PR context pre-loaded.
  \item[\code{/rename "feature-auth"}] Give descriptive names for later retrieval.
\end{description}

\subsection{The CURRENT\_WORK.md Pattern}

\practitioner{Maintain a \code{CURRENT\_WORK.md} file at project root.
It provides 30-second context resume after any interruption.}

\begin{minted}{markdown}
## Right Now
Refactoring the payment processing module to use async handlers.

## Why
Current sync handlers cause 200ms latency spikes under load.

## Next Step
Write integration tests for the new async PaymentProcessor class.

## Context When I Return
- PaymentProcessor class is in src/payments/processor.py
- 3 of 5 handler methods converted, 2 remaining
- Tests in tests/payments/ — existing tests still pass
\end{minted}

\margintip{Update CURRENT\_WORK.md at the start and end of every session. 60 seconds invested saves 10 minutes of re-discovery.}

This pattern is especially valuable when:
\begin{itemize}
  \item You work on a project intermittently (days between sessions)
  \item Multiple people work on the same project
  \item You context-switch frequently between projects
\end{itemize}

% -------------------------------------------------------------------
\section{Auto-Memory and Cross-Session Persistence}
\label{sec:memory}

\marginnote[Official]{Claude records and recalls memories across sessions automatically. Memory files persist in \code{\textasciitilde/.claude/projects/*/memory/}.}

\official{Claude maintains auto-memory that persists across conversations.}
It records patterns, decisions, and lessons learned, making them available
in future sessions without explicit re-prompting.

\subsection{What to Store in Memory}

\begin{description}
  \item[Stable patterns] Confirmed conventions that apply across sessions
  \item[Architecture decisions] Why you chose X over Y (not just what)
  \item[Hard-won lessons] Bugs, workarounds, gotchas that took time to discover
  \item[Quality metrics] Baselines and targets for dashboards
  \item[Directory conventions] Non-obvious file organization patterns
  \item[Known issues] Pre-existing problems to avoid re-investigating
\end{description}

\subsection{What NOT to Store}

\begin{description}
  \item[Transient state] Current task details, in-progress work
  \item[Unverified conclusions] Patterns seen once, not confirmed
  \item[Duplicates of CLAUDE.md] Memory supplements configuration, does not replace it
  \item[Speculative notes] Only store what you have verified
\end{description}

\practitioner{The test for whether something belongs in memory:
``Will this still be true and useful in three months?''
If yes, store it. If uncertain, wait for confirmation.}

% -------------------------------------------------------------------
\section{Plan Documents for Large Tasks}
\label{sec:plans}

\marginnote[Official]{Use Plan Mode (Shift+Tab) for read-only analysis before implementation.\sourceurl{https://code.claude.com/docs/en/common-workflows}}

\official{The recommended workflow for complex tasks:
Explore $\rightarrow$ Plan $\rightarrow$ Implement $\rightarrow$ Commit.}

Plan Mode lets Claude analyze your codebase without modifying anything.
Use it for architecture discussions, impact analysis, and strategy decisions.

\subsection{Plan Document Lifecycle}

\practitioner{For tasks exceeding one hour or 500 lines of changes,
create a plan document with lifecycle management:}

\begin{enumerate}
  \item \textbf{Active} --- currently being planned, open questions remain
  \item \textbf{Ready} --- plan approved, implementation can begin
  \item \textbf{Implemented} --- all phases complete
  \item \textbf{Archived} --- moved to \code{docs/plans/archived/}
\end{enumerate}

\subsection{The ``Decisions Made'' Section}

\practitioner{Every plan document should include a ``Decisions Made'' section
documenting:}

\begin{itemize}
  \item Why specific implementations were chosen
  \item Alternatives considered but rejected (and why)
  \item Assumptions made
  \item Tradeoffs accepted (performance, complexity, scope)
\end{itemize}

This section is invaluable for postmortem analysis.
When a decision turns out to be wrong, you can trace back to the reasoning
and understand what information was missing.

\begin{keyinsight}
Context management is not about minimizing Claude's input---it is about
maximizing the signal-to-noise ratio. Every mechanism in this chapter
(clear, compact, sessions, memory, plans) serves the same purpose:
ensuring that Claude's limited attention is spent on what matters most
for the current task.
\end{keyinsight}
